Este apêndice registra onde estão os dados usados para sustentar a análise
quantitativa do Capítulo~\ref{chapter:resultados-testes}. A campanha
\texttt{full\_run} corresponde à última rodada completa de testes executada com a
versão do firmware analisada nesta monografia. No repositório local, os artefatos
estão em \path{transmitter/testing/results/full_run}. A mesma campanha pode ser
consultada no
\href{https://github.com/cirillom/WCAN/tree/3b57f282a41413706ea0240b9030d72ff5438133/transmitter/testing/results/full_run}{repositório GitHub}.

A pasta contém o plano de execução, a configuração das placas, a configuração das
suítes, o resumo produzido pelo runner, o relatório consolidado do analisador e os
logs seriais de cada placa. Cada execução possui um \texttt{scenario.json} e um
conjunto de logs separados por papel e porta serial. Essa organização permite
recalcular as métricas agregadas e inspecionar manualmente casos reprovados.

Na matriz nominal, a campanha possuía 1116 execuções. O analisador consolidou 1115
pastas com logs suficientes e classificou 847 execuções como aprovadas. Como uma
reprovação pode envolver entrega, transporte, cenário, heap, crash ou frequência
observada, os resultados foram interpretados por classe de falha, não apenas pela
taxa bruta de aprovação.
